% English translation of checked_pages/esq_119_028.tex; source: cote 119, site 28, PDF page 29, typescript page 26.
% AI-translated by GPT-6 Astra (OpenAI, ChatGPT Pro mode).
\scanpage{119/28}
\mspage{26}

The compactified Mumford-Deligne moduli multiplicities $M^{\wedge}_{g,\nu}$
for stable algebraic curves of type $(g,\nu)$
provide a typical and particularly inspiring example, which
played an important motivating role in the resumption of my reflection
on stratified structures, from December 1981 to January
1982. Two-dimensional geometry provides many other
examples of such stratified moduli structures, all of which,
moreover (apart from expedients of rigidification\frapparatus{V-FR-028-001}), appear
as ``multiplicities'' rather than as spaces or varieties
in the ordinary sense (the points of these multiplicities\frapparatus{V-FR-028-002} can have
nontrivial automorphism groups). Among the objects of
two-dimensional geometry giving rise to such stratified moduli
structures of arbitrary dimension, or even infinite
dimension, I shall mention polygons (euclidean, spherical or
hyperbolic), systems of lines in a plane (projective,
say), systems of ``pseudolines'' in a topological projective
plane, or more general immersed curves with normal \unsure{crossings}\tnote{TN-028-01}{The French alternatives are ``crc isements|croisements''. Both are represented by ``crossings'', with the unresolved letterform still marked and linked.}
in a given surface (compact, say).

The simplest nontrivial example of a stratified
structure is obtained by considering a pair $(X,Y)$ consisting of a space
$X$ and a closed subspace $Y$, making a suitable equisingularity
assumption on $X$ along $Y$, and assuming
further (to fix ideas) that the two strata $Y$ and $X\setminus Y$ are
topological \uline{manifolds}. The naive idea, in such a situation,
is to take ``the'' tubular neighborhood $T$ of $Y$ in $X$, whose
boundary $\partial T$ should also be a smooth manifold, fibered over
$Y$ with smooth compact fibers, $T$ itself being identified with
the bundle of cones over $T$ associated with the preceding bundle. Setting
\[ U = X \setminus \mathrm{Int}(T), \]
one finds a manifold with boundary whose boundary is canonically
isomorphic to that of $T$. That said, the envisaged ``building blocks''
are the manifold with boundary $U$ (compact when $X$ was compact, and
replacing the ``open'' stratum $X\setminus Y$ while making it more precise) and the manifold

\margin{28}

